% Unite: papier 34, section 24; traduction francaise depuis la source allemande auditee.

\providecommand{\mo}{\mathfrak{o}}
\providecommand{\mc}{\mathfrak{c}}
\providecommand{\mM}{\mathfrak{M}}
\providecommand{\mR}{\mathfrak{R}}

\subsection*{\S{} 24. Traces et caractères}
Si, dans une représentation $\mo\sim\mathfrak D$ d'un système hypercomplexe $\mo$, la matrice $A$ est associée à l'élément $a$, on pose
\[
        \operatorname{Sp}_{\mathfrak D}(a)=\operatorname{Sp} A.
\]
Les traces sont des fonctions linéaires:
\[
        \operatorname{Sp}(c+d)=\operatorname{Sp}(c)+\operatorname{Sp}(d),
        \qquad
        \operatorname{Sp}(c\alpha)=\alpha\operatorname{Sp}(c).
\]
Les représentations équivalentes ont les mêmes traces. La trace dans une représentation réductible est la somme des traces dans les représentations induites par les facteurs de composition.

\emph{Démonstration.} En effet,
\[
        \operatorname{Sp}\begin{pmatrix}
        A_{11}&0&\cdots&0\\
        *&A_{22}&\cdots&0\\
        \vdots&\vdots&\ddots&\vdots\\
        *&*&\cdots&A_{ss}
        \end{pmatrix}
        =\operatorname{Sp}(A_{11})+\operatorname{Sp}(A_{22})+\cdots+\operatorname{Sp}(A_{ss}).
\]
On appelle trace principale la trace dans la représentation régulière. On appelle trace réduite la somme des traces dans les différentes représentations irréductibles.

Si $c$ est un élément de l'idéal nilpotent maximal $\mc$, alors $\operatorname{Sp}(c)=0$ pour toute représentation.

\emph{Démonstration.} Il suffit d'étudier les représentations définies par les idéaux à gauche simples de $\mo/\mc$; en effet, dans toute autre représentation, seuls les facteurs de composition ainsi obtenus apparaissent, et la trace se compose additivement. Or $c$ annule tous les éléments de $\mo/\mc$; à tout $c\in\mc$ est donc associée la matrice nulle, et par conséquent $\operatorname{Sp}(c)=0$.

Si $\mo$ est simple bilatère et si $P$ est algébriquement clos, donc si $\mo=\sum c_{ik}P$ est un anneau de matrices, alors, dans la représentation irréductible de $\mo$, à l'élément $a=\sum c_{ik}\alpha_{ik}$ est associée la matrice $(\alpha_{ik})$. Donc
\[
        \operatorname{Sp}_{\rm red}(a)=\sum_i\alpha_{ii},
        \qquad
        \operatorname{Sp}_{\rm pr}(a)=n\,\operatorname{Sp}_{\rm red}(a)=n\sum_i\alpha_{ii}.
\]
En particulier
\[
        \operatorname{Sp}_{\rm red}(c_{ik})=0\ (i\ne k),\quad
        \operatorname{Sp}_{\rm red}(c_{ii})=e,
\]
et
\[
        \operatorname{Sp}_{\rm pr}(c_{ik})=0\ (i\ne k),\quad
        \operatorname{Sp}_{\rm pr}(c_{ii})=n e.
\]

Lorsqu'on passe de $P$ au corps algébriquement clos $\Omega$, la trace principale ne change pas, puisqu'elle peut être calculée à partir de la même base. Il n'en va pas de même de la trace réduite.

Les traces des éléments $a$ du système sans radical, dans les représentations absolument irréductibles, s'appellent les caractères; on les désigne par $\chi(a)$ ou, lorsqu'il faut indiquer de quelle représentation il s'agit, par $\chi^{(\nu)}(a)$.

Dans une représentation irréductible de degré $n_\nu$, les éléments du centre sont représentés, d'après le \S{}21, par des matrices diagonales $E\theta^{(\nu)}$, où $z\mapsto\theta^{(\nu)}(z)$ est une représentation du premier degré du centre (ou un homomorphisme du centre dans $\Omega$). Il en résulte pour la trace la valeur $n_\nu\theta^{(\nu)}(z)$. Donc les homomorphismes du centre sont liés aux caractères par
\[
        \chi^{(\nu)}(z)=n_\nu\theta^{(\nu)}(z)                 \tag{1}
\].
Dans le cas commutatif, $n_\nu=1$, et les caractères eux-mêmes donnent les homomorphismes (voir \S{}22). Si le corps $\Omega$ est de caractéristique zéro, ce que nous supposerons toujours dans la suite, on peut diviser (1) par $n_\nu$:
\[
        \theta^{(\nu)}(z)=\frac{\chi^{(\nu)}(z)}{n_\nu}.
\]
La propriété d'homomorphisme s'exprime par
\[
        \frac{\chi^{(\nu)}(z+z')}{n_\nu}
        =\frac{\chi^{(\nu)}(z)}{n_\nu}+\frac{\chi^{(\nu)}(z')}{n_\nu},
\]
\[
        \frac{\chi^{(\nu)}(zz')}{n_\nu}
        =\frac{\chi^{(\nu)}(z)}{n_\nu}\frac{\chi^{(\nu)}(z')}{n_\nu}
\].

Une classe de représentation est déjà déterminée de façon univoque par les seules traces des matrices. (Pour connaître les traces de toutes les matrices, il suffit naturellement de connaître les traces des éléments de base dans la représentation donnée.)

\emph{Démonstration.} Le module de représentation $\mR$ est connu lorsqu'on sait combien de fois chaque module de représentation irréductible $\mM_\nu$ apparaît comme facteur direct. Si ce nombre est $\mu_\nu$, la trace de $e^{(\nu)}$ est égale à $\mu_\nu n_\nu$, donc
\[
        \mu_\nu=\frac{\operatorname{Sp}(e^{(\nu)})}{n_\nu}.
\]
